\begin{abstract}
A detector can retain a low false-positive rate after a change of source while missing most text from a particular generator. Conversely, high sensitivity can coexist with widespread false attribution of human text. We examine these distinctions in a completed, locally executed study of seven frozen detectors, three source domains, and three open-weight model configurations. For 450 source items, we generate three direct passages, all nine ordered generator--rewriter combinations, and three rewrites of the associated human excerpt. A separate cohort supplies 1,500 human controls. The resulting 6,750 generations and 8,700 texts yield 60,900 detector scores without retraining or threshold adaptation. Binoculars and Fast-DetectGPT flag 42/1,500 and 37/1,500 human controls, respectively, but their direct-generation sensitivity varies markedly by model. RADAR flags no human controls and only 63/1,350 direct generations. Gemma-to-Qwen rewriting reduces Binoculars flags from 255/450 to 17/450, yet Gemma self-rewriting also reduces them to 44/450. Length-contract failures affect 1,483/6,750 generations and remain in every planned denominator. We report all detector and history cells, reproduce the fixed 357-cell uncertainty family, and distinguish descriptive score transport from meaning-preserving evasion. The latter remains pending independent blinded review. The study supports evaluating specificity, production history, same-model controls, and completion constraints together rather than interpreting a direct-to-rewrite detection drop in isolation.
\end{abstract}
